\section{Theoretical Overview}\label{sec:theoretical_overview} % (fold)
Based on the project specification, the RISC CPU designed in this assignment features a 3-bit opcode and a 5-bit operand. This architecture supports 8 types of instructions (HLT, SKZ, ADD, AND, XOR, LDA, STO, JMP) and a 32-word memory address space. The processor operates synchronously with a clock signal and an active-high reset signal (the default reset state is \texttt{INST\_ADDR}). Program execution stops upon encountering the \texttt{HLT} instruction.

By analyzing the controller state table and the functional requirements of each block, the CPU's operation flow can be divided into two primary phases spanning 8 clock cycles: the Fetch Phase and the Execution Phase.

% TODO: Insert table content:
\begin{table}[H]
	\caption{CPU Controller state and Control signals}
	% \label{tab:controller_states}
	\begin{center}
		\begin{tabular}[c]{l|l}
			\hline
			\multicolumn{1}{c|}{\textbf{}} &
			\multicolumn{1}{c}{\textbf{}}      \\
			\hline
			a                              & b \\
			c                              & d \\

			\hline
		\end{tabular}
	\end{center}
\end{table}


\subsection{Fetch Phase}
\begin{itemize}
	\item \textbf{State 1: INST\_ADDR.} The Program Counter (PC) provides the address of the current instruction. The Address Mux selects the PC address (\texttt{sel = 1}) to access the memory.
	\item \textbf{State 2: INST\_FETCH.} The memory read signal (\texttt{rd = 1}) is asserted. The instruction data is retrieved from Memory.
	\item \textbf{State 3: INST\_LOAD.} The fetched instruction is loaded into the Instruction Register (IR) via the \texttt{ld\_ir = 1} control signal.
	\item \textbf{State 4: IDLE.} An intermediate wait state for the system to stabilize.
\end{itemize}

\subsection{Execution Phase}
\begin{itemize}
	\item \textbf{State 5: OP\_ADDR.} The IR splits the instruction into a 3-bit Opcode and a 5-bit Operand Address. The Address Mux switches to select the operand address from the IR (\texttt{sel = 0}). Simultaneously, the PC increments (\texttt{inc\_pc = 1}) to prepare for the next instruction cycle.
	\item \textbf{State 6: OP\_FETCH.} If the instruction requires data from memory (e.g., ADD, AND, LDA), the CPU reads the data from the operand address in the Memory (\texttt{rd = 1}).
	\item \textbf{State 7: ALU\_OP.} The Arithmetic Logic Unit (ALU) performs the required computation between the data fetched from Memory and the data currently held in the Accumulator (AC), based on the decoded Opcode. If the instruction is \texttt{SKZ} and the \texttt{zero} flag condition is met, the PC increments once more (\texttt{inc\_pc = 1}) to skip the subsequent instruction.
	\item \textbf{State 8: STORE.} The result of the operation is either stored back into the Memory (for the \texttt{STO} instruction, asserting \texttt{wr = 1} and \texttt{data\_e = 1}) or loaded into the Accumulator (for arithmetic and logic instructions, asserting \texttt{ld\_ac = 1}).
\end{itemize}

% section Theoretical Overview (end)
